\documentclass[12pt,a4paper]{article}

% ----- Encoding and language -----------------------------------------
\usepackage[utf8]{inputenc}
\usepackage[T1]{fontenc}
\usepackage{lmodern}
\usepackage[spanish,es-tabla,es-nodecimaldot, es-noquoting]{babel}
\usepackage{microtype}

% ----- Page layout --------------------------------------------------- 
\usepackage[a4paper,left=2.6cm,right=2.6cm,top=2.8cm,bottom=2.6cm,headheight=15pt]{geometry}
\linespread{1.15}
\setlength{\parskip}{4pt}
\setlength{\parindent}{0pt}

% ----- Colors --------------------------------------------------------
\usepackage[table]{xcolor}
\definecolor{darkblue}{HTML}{0B2545}
\definecolor{mediumblue}{HTML}{1D4E89}
\definecolor{lightblue}{HTML}{DCE7F3}
\definecolor{borderblue}{HTML}{B8C7DD}
\definecolor{lightgray}{HTML}{F5F6F8}

% ----- Packages ------------------------------------------------------
\usepackage{graphicx,float,booktabs,tabularx,array,multirow,longtable}
\usepackage{enumitem,amsmath,amssymb,caption,xparse,url}
\usepackage{siunitx}
\sisetup{output-decimal-marker={,}, inter-unit-product=\ensuremath{{}\cdot{}}}
\usepackage{tikz}
\usetikzlibrary{arrows.meta,calc,positioning,patterns}
\usepackage{pgfplots}
\pgfplotsset{compat=1.16}
\usepackage[most]{tcolorbox}
\usepackage{fancyhdr,titlesec}

% ----- Header and footer ---------------------------------------------
\pagestyle{fancy}
\fancyhf{}
\fancyfoot[L]{\sffamily\footnotesize\bfseries\color{darkblue} Facultad de
Ingeniería Industrial}
\fancyfoot[R]{\sffamily\footnotesize\bfseries\color{darkblue}Pag \thepage}
\renewcommand{\headrulewidth}{0pt}
\renewcommand{\footrulewidth}{0.8pt}
\renewcommand{\footrule}{{\color{darkblue}\hrule width\headwidth
height\footrulewidth}}

% ----- Section titles -----------------------------------------------
\titleformat{\section}
  {\sffamily\bfseries\Large\color{darkblue}}{\thesection.}{0.6em}{}
  [{\vspace{-0.6ex}\color{mediumblue}\titlerule[1.3pt]}]
\titleformat{\subsection}
  {\sffamily\bfseries\large\color{mediumblue}}{\thesubsection}{0.6em}{}
\titleformat{\subsubsection}
  {\sffamily\bfseries\normalsize\color{darkblue}}{\thesubsubsection}{0.6em}{}
\titlespacing*{\section}{0pt}{2.4ex plus 1ex minus .2ex}{1.6ex}
\titlespacing*{\subsection}{0pt}{2.0ex plus .8ex minus .2ex}{1.0ex}
\titlespacing*{\subsubsection}{0pt}{1.6ex plus .5ex}{0.6ex}
\addto\captionsspanish{\renewcommand{\contentsname}{CONTENIDO}}

% ----- Captions and tables ------------------------------------------
\captionsetup{font=small,labelfont={bf,color=darkblue},labelsep=period,
              textfont={color=black!80},skip=6pt}
\renewcommand{\arraystretch}{1.22}
\setlist{itemsep=2pt,topsep=3pt,parsep=0pt}

% Table header cell (white text on dark blue)
\newcommand{\encab}[1]{\textcolor{white}{\sffamily\bfseries #1}}

% ----- Text boxes ----------------------------------------------------
\newtcolorbox{nota}[1][Nota]{colback=lightblue,colframe=mediumblue,coltitle=white,
  fonttitle=\sffamily\bfseries\small,title=#1,arc=2pt,boxrule=0.6pt,
  left=7pt,right=7pt,top=4pt,bottom=4pt,breakable}
\newtcolorbox{enunciado}{colback=lightblue,colframe=darkblue,arc=2pt,boxrule=0pt,
  leftrule=3pt,left=8pt,right=8pt,top=4pt,bottom=4pt,breakable,fontupper=\sffamily\small}

% ----- Image with placeholder box if file does not exist -------------
% Usage \image[width][box height]{file}{description}
\NewDocumentCommand{\image}{O{0.6} O{5cm} m m}{%
  \IfFileExists{#3}{\includegraphics[width=#1\linewidth]{#3}}{%
    {\setlength{\fboxsep}{0pt}\fcolorbox{borderblue}{lightgray}{%
      \parbox[c][#2][c]{#1\linewidth}{\centering\sffamily\small\color{darkblue}%
      \textbf{[ Insert image
]}\\[3pt]\texttt{\detokenize{#3}}\\[3pt]\footnotesize #4}}}}}

% ----- Links ---------------------------------------------------------
\usepackage[colorlinks=true, breaklinks=true, linkcolor=darkblue,
urlcolor=mediumblue, citecolor=darkblue,
  pdftitle={Fabricación de un buje de bronce con ajuste H7/h6},
  pdfauthor={Integrantes del grupo},
  pdfsubject={Procesos de Manufactura — Trabajo de investigación —
UNMSM}]{hyperref}
\urlstyle{same}
\makeatletter\g@addto@macro\UrlBreaks{\do\-}\makeatother
\Urlmuskip=0mu plus 2mu\relax
\setlength{\emergencystretch}{2.5em}
\newcommand{\refbib}[1]{[\ref{#1}]}

% =====================================================================

% ----- NUEVO: unidades y atajos para los cálculos -------------------
\DeclareSIUnit{\rev}{rev}
\DeclareSIUnit{\rpm}{rpm}
\newcommand{\umm}{\,\mathrm{mm}}
\newcommand{\umin}{\,\mathrm{min}}
\newcommand{\urpm}{\,\mathrm{rpm}}
\newcommand{\umpm}{\,\mathrm{m/min}}
\newcommand{\ummrev}{\,\mathrm{mm/rev}}
\newcommand{\ummmin}{\,\mathrm{mm/min}}
\newcommand{\ucmmin}{\,\mathrm{cm^3/min}}
\newcommand{\uum}{\,\mu\mathrm{m}}

% =====================================================================
\begin{document}
% =====================================================================

% ===== COVER PAGE ====================================================
\begin{titlepage}
	\thispagestyle{empty}
	\begin{tikzpicture}[remember picture,overlay]
		\fill[darkblue] (current page.north west) rectangle ($(current page.north
			east)+(0,-4.4cm)$);
		\fill[mediumblue] ($(current page.north west)+(0,-4.4cm)$) rectangle
		($(current page.north east)+(0,-4.65cm)$);
		\fill[darkblue] (current page.south west) rectangle ($(current page.south
			east)+(0,1.1cm)$);
		\node[text=white,align=center,text width=16cm,font=\sffamily] at ($(current
			page.north)+(0,-2.25cm)$) {%
			{\Large\bfseries UNIVERSIDAD NACIONAL MAYOR DE SAN MARCOS}\\[5pt]
			{\small\itshape Universidad del Perú, Decana de América}\\[10pt]
			{\normalsize\bfseries ESCUELA PROFESIONAL DE INGENIERÍA INDUSTRIAL}};
	\end{tikzpicture}

	\vspace*{2cm}

	\begin{center}
		\image[0.22][4.3cm]{img/unmsm-logo.png}{}

		\vspace{0.8cm}
		{\sffamily\bfseries\color{darkblue}\fontsize{21}{25}\selectfont
			Fabricación de un buje de bronce\\con ajuste H7/h6\par}

		\vspace{0.3cm}
		{\sffamily\large\color{mediumblue} Trabajo de investigación: proceso de
			torneado\par}

		\vspace{0.4cm}

		\begin{tcolorbox}[colback=white,colframe=white,arc=3pt,width=0.9\linewidth,
				box align=center,left=12pt,right=12pt,top=8pt,bottom=8pt]
			\centering\sffamily
			\begin{tabular}{@{}>{\bfseries\color{darkblue}}r@{\hspace{1em}}l@{}}
				Asignatura: & Procesos de Manufactura (L1-3) \\[5pt]
				Docente:    & Mg. Ing. Víctor Rosales Urbano \\
			\end{tabular}
		\end{tcolorbox}

		\vspace{0.4cm}
		{\sffamily\bfseries\color{darkblue}\large Integrantes\par}
		\vspace{5pt}
		\sffamily\begin{tabular}{@{}l@{\hspace{2.5em}}r@{}}
			Apellidos, Nombres del integrante 1 & (código) \\
			Apellidos, Nombres del integrante 2 & (código) \\
			Apellidos, Nombres del integrante 3 & (código) \\
			Apellidos, Nombres del integrante 4 & (código) \\
		\end{tabular}
	\end{center}

	\vfill
	\begin{center}
		{\sffamily\bfseries\color{darkblue}\large Lima, octubre de 2026\par}
	\end{center}
\end{titlepage}
\setcounter{page}{2}

% ===== CONTENIDO =====================================================
\tableofcontents
\newpage

% ===== 1. INTRODUCCIÓN ===============================================
\section{INTRODUCCIÓN}

El buje (o casquillo) es uno de los elementos de máquina más sencillos y, a la
vez, más frecuentes: un cilindro hueco que sirve de superficie de apoyo y guía
para un eje. Su función depende casi por completo de dos diámetros y de la
relación entre ellos, es decir, del \emph{ajuste}. Un buje con demasiado juego
deja de guiar; uno con muy poco juego se agarrota. Por eso su fabricación es un
buen ejemplo de cómo el proceso de manufactura (torneado) y la tolerancia
especificada en el plano se condicionan mutuamente.\\[-4pt]

Este trabajo de investigación plantea, de forma \textbf{teórica}, la
fabricación de un buje de bronce con ajuste \textbf{H7/h6} mediante torneado.
Se parte de las fórmulas de mecanizado vistas en el curso (velocidad de corte,
avance, caudal de viruta, número de pasadas y tiempos) y se aplican paso a paso
a cada operación, de modo que el documento pueda leerse como una hoja de
proceso completa y justificada.

\subsection{Objetivo del proyecto}

\textbf{Objetivo general.} Planificar la fabricación por torneado de un buje de
bronce SAE 660 con ajuste H7/h6, determinando la secuencia de operaciones, los
parámetros de corte y los tiempos de mecanizado, y verificando analíticamente
que el proceso es capaz de cumplir las tolerancias especificadas.

\textbf{Objetivos específicos.}
\begin{itemize}
	\item Definir la geometría de la pieza y las tolerancias dimensionales,
	      geométricas y de acabado superficial que exige el ajuste H7/h6.
	\item Seleccionar el material de partida, la máquina y las herramientas de
	      corte adecuadas para un bronce al plomo.
	\item Calcular, con las fórmulas del curso, la velocidad de corte, las rpm,
	      el avance, la profundidad y el número de pasadas, el caudal de viruta y
	      los tiempos de cada operación.
	\item Elaborar la hoja de ruta y el plano de fabricación de la pieza.
	\item Estimar el presupuesto de errores del proceso y compararlo con la banda
	      de tolerancia, proponiendo mejoras.
\end{itemize}

\subsection{Descripción de la pieza}

La pieza es un buje cilíndrico de pared delgada, sin brida, con las
características de la Tabla~\ref{tab:pieza}. Su agujero ($\varnothing 25$ H7)
aloja un eje $\varnothing 25$ h6, y su diámetro exterior ($\varnothing 32$ p6)
se monta a presión en un alojamiento H7, de modo que el buje queda fijo en la
carcasa y el eje gira (o se desliza) dentro de él.

\begin{table}[!htb]
	\centering
	\caption{Especificaciones generales del buje}\label{tab:pieza}
	\small
	\rowcolors{1}{white}{lightblue}
	\begin{tabularx}{\linewidth}{>{\bfseries}l X}
		Denominación          & Buje cilíndrico de bronce (sin brida)
		\\
		Material              & Bronce al plomo SAE 660 (UNS C93200), barra
		redonda $\varnothing 35$ mm                                                                                                       \\
		Diámetro interior     & $\varnothing 25$ H7 $=
			                                            25{,}000^{+0{,}021}_{\;\;0}$ mm                                                      \\
		Diámetro exterior     & $\varnothing 32$ p6 $= 32^{+0{,}042}_{+0{,}026}$
		mm                                                                                                                                \\
		Longitud              & $40 \pm 0{,}1$ mm
		\\
		Chaflanes             & $1\times45^\circ$ en el exterior y
		$0{,}5\times45^\circ$ en el interior                                                                                              \\
		Acabado superficial   & Ra $\le 0{,}8$ \textmu m en el agujero; Ra $\le
			                                                                   1{,}6$ \textmu m en el exterior                               \\
		Tolerancia geométrica & Coaxialidad agujero--exterior $\le \varnothing 0{,}02$ mm; perpendicularidad de las caras $\le 0{,}02$ mm \\
		Tolerancia general    & ISO 2768-m para cotas sin tolerancia indicada                                                             \\
		Masa aproximada       & 0,11 kg ($V\approx 12{,}5$ cm$^3$)                                                                        \\
	\end{tabularx}
\end{table}

\textbf{¿Por qué H7/h6?} En el sistema ISO de ajustes, H7/h6 es un ajuste de
\emph{juego mínimo} (deslizante o de posicionamiento): el agujero H7 tiene su
desviación inferior en cero y el eje h6 tiene su desviación superior en cero,
de modo que en el peor caso el juego es nulo y en el mejor caso es solo de unas
pocas centésimas de milímetro. Es adecuado para guiar con precisión piezas que
giran o se deslizan lentamente y con lubricación, como casquillos guía,
pasadores y ejes de mecanismos de precisión; no lo es para ejes que giran a alta
velocidad, donde se requiere un juego mayor para formar la película de aceite
(por ejemplo, H7/f7).

\begin{nota}[Alcance y supuestos del trabajo]
	No se fabricó la pieza: todos los valores son resultado de cálculo.
	Las velocidades de corte se tomaron de la tabla del curso (Anexo~A); los
	valores de propiedades del material son típicos de la literatura (deben
	confirmarse con el certificado del proveedor); y los tiempos auxiliares
	(carga, cambio de herramienta, inspección) son estimaciones razonadas. Se
	asume un torno CNC de bancada inclinada con velocidad de husillo máxima de
	4\,000 rpm.
\end{nota}

% ===== 2. MARCO TEÓRICO ==============================================
\section{MARCO TEÓRICO}

\subsection{El torno y su selección}

El \textbf{torno} es la máquina-herramienta que fabrica \emph{piezas de
	revolución}: la pieza gira sujeta en el plato del husillo y una herramienta de
un solo filo, montada en la torreta o en el carro, avanza contra ella y arranca
viruta. Sus partes principales (Figura~\ref{fig:partes}) son el \emph{plato} con
sus \emph{mordazas}, que sujetan y centran la pieza; el \emph{cabezal} o husillo
principal, que le da el giro; la \emph{torreta}, que porta las herramientas; los
\emph{carros} longitudinal y transversal, que producen el avance; y la
\emph{bancada}, que soporta y guía todo el conjunto. En el torno paralelo
convencional se añaden la caja de velocidades, el tornillo patrón y la
contrapunta; en el torno CNC estos movimientos los gobierna el control numérico.

\begin{figure}[H]
	\centering
	\image[0.62][5.2cm]{img/partes_torno_cnc.png}{Partes de un torno CNC: plato, mordazas y torreta}
	\caption{Partes principales de un torno CNC (imagen referencial).}\label{fig:partes}
\end{figure}

La Tabla~\ref{tab:tornos} resume los tipos de torno y su idoneidad para fabricar
un buje.

\begin{table}[!htb]
	\centering
	\caption{Tipos de torno y su idoneidad para el buje}\label{tab:tornos}
	\small
	\rowcolors{2}{white}{lightblue}
	\begin{tabularx}{\linewidth}{l X c}
		\rowcolor{darkblue}
		\encab{Tipo}                                    & \encab{Característica}
		                                                & \encab{Idoneidad}
		\\
		Paralelo                                        & Máquina universal, de operación manual; hoy se
		reserva a trabajos puntuales o de mantenimiento & Baja                                                                                             \\
		Copiador                                        & Reproduce una réplica de una plantilla o guía
		                                                & Baja                                                                                             \\
		Revólver                                        & Varias herramientas trabajan casi simultáneamente para
		reducir tiempo en series                        & Media                                                                                            \\
		Automático                                      & Alimentación de barra y ciclo automático, para grandes
		series                                          & Media
		\\
		Vertical                                        & Eje vertical; fabrica piezas de gran tamaño
		                                                & Nula                                                                                             \\
		\textbf{CNC}                                    & Dirigido por computadora; gran capacidad de producción y precisión de mecanizado & \textbf{Alta} \\
	\end{tabularx}
\end{table}

Para un buje con tolerancia IT6--IT7 y dos diámetros que deben quedar coaxiales,
conviene el \textbf{torno CNC}: permite hacer refrentado, taladrado, mandrinado,
cilindrado y tronzado en una sola sujeción de la barra, lo que minimiza el error
de coaxialidad, y repite las mismas condiciones de corte de pieza a pieza.

\subsection{Operaciones de torneado}

El torneado agrupa varias operaciones según la trayectoria de la herramienta
respecto del eje de giro (Tabla~\ref{tab:operaciones}). En el buje intervienen
seis de ellas.

\begin{table}[!htb]
	\centering
	\caption{Operaciones de torneado y su uso en el buje}\label{tab:operaciones}
	\small
	\rowcolors{2}{white}{lightblue}
	\begin{tabularx}{\linewidth}{l X c}
		\rowcolor{darkblue}
		\encab{Operación}                & \encab{Descripción}
		                                 & \encab{¿Se usa?}
		\\
		Cilindrado exterior              & Da forma exterior y reduce el diámetro de
		la barra                         & Sí                                                                           \\
		Cilindrado interior (mandrinado) & Cilindros o conos interiores y ranuras interiores en piezas agujereadas & Sí \\
		Refrentado                       & Superficie plana perpendicular al eje de giro                           & Sí \\
		Tronzado                         & Seccionamiento de la barra o pieza con una herramienta de tronzar       & Sí \\
		Forma / chaflanado               & Perfiles definidos por el filo o la trayectoria; incluye los chaflanes  & Sí \\
		Taladrado (en torno)             & Agujero axial con broca montada en el
		contrapunto o la torreta         & Sí
		\\
		Roscado                          & Varias pasadas a lo largo de la longitud de la rosca                    & No \\
		Ranurado                         & Abre ranuras o canales en las piezas                                    & No \\
		Moleteado                        & Produce una superficie áspera o rugosa con moletas                      & No \\
		Cónico                           & Superficies cónicas por inclinación del carro o interpolación           & No \\
	\end{tabularx}
\end{table}

En toda operación de arranque conviene distinguir dos tipos de pasada
(Tabla~\ref{tab:desbaste}). El \textbf{desbaste} quita la mayor cantidad de
material para acercarse a la medida final; el \textbf{acabado} es una pasada casi
superficial que da la tolerancia y la rugosidad. Se desbasta con velocidad de
corte moderada, avance alto y gran profundidad; se acaba con velocidad alta,
avance bajo y poca profundidad.

\begin{table}[!htb]
	\centering
	\caption{Desbaste frente a acabado}\label{tab:desbaste}
	\small
	\rowcolors{2}{white}{lightblue}
	\begin{tabular}{l c c}
		\rowcolor{darkblue}
		\encab{Parámetro}    & \encab{Desbaste} & \encab{Acabado} \\
		Velocidad de corte   & baja             & alta            \\
		Avance               & alto             & bajo            \\
		Profundidad de corte & grande           & pequeña         \\
		Tiempo de mecanizado & bajo             & alto            \\
	\end{tabular}
\end{table}

\subsection{Fórmulas de mecanizado}

Las siguientes relaciones, vistas en clase~\refbib{bib:clase}, se emplean en los
cálculos de la sección~\ref{sec:calculos}. La nomenclatura completa se resume en
el Anexo~B.\\[-4pt]

\textbf{Velocidad de corte.} Es la velocidad tangencial en la superficie de la
pieza. Con $n$ en rev/min y $1\ \text{rev}=2\pi$:
\begin{equation}\label{eq:v}
	v = n\left(\frac{D_i}{2}\right)\ \Longrightarrow\ v=\frac{\pi\, D_i\, n}{1000}
	\quad[\text{m/min};\ D_i\ \text{en mm};\ n\ \text{en rpm}]
\end{equation}
Despejando, la velocidad de giro que corresponde a una velocidad de corte
recomendada es $n = 1000\,v/(\pi D_i)$.

\textbf{Velocidad lineal o de avance.}
\begin{equation}\label{eq:vl}
	v_L = n\cdot a \qquad [\text{mm/min}]
\end{equation}
donde $a$ es el avance por revolución (mm/rev).

\textbf{Caudal de viruta.} Volumen de material arrancado por unidad de tiempo:
\begin{equation}\label{eq:Q}
	Q = v_L\cdot A
\end{equation}
con $A$ el área de la corona circular removida en la pasada,
$A=\frac{\pi}{4}\left(D_i^2-D_f^2\right)$.

\textbf{Profundidad de corte y número de pasadas.} La profundidad total radial
que debe quitarse se reparte entre pasadas de desbaste y de acabado, y el número
de pasadas debe ser entero:
\begin{equation}\label{eq:pT}
	p_T=\frac{D_i-D_f}{2}=N_D\,P_D+N_A\,P_A
\end{equation}
Valores usuales: acabado $P_A\approx 0{,}1$--$0{,}5$ mm; desbaste $P_D\approx
	0{,}4$--$3$ mm.

\textbf{Tiempo de pasada y tiempo de mecanizado.}
\begin{equation}\label{eq:tp}
	t_p=\frac{L+y+\Delta}{v_L}=\frac{L+y+\Delta}{a\cdot n},
	\qquad t_m = N\left(\frac{L+y+\Delta}{a\cdot n}\right)
\end{equation}
donde $L$ es la longitud mecanizada, $y$ el retiro (aproximación) de la
herramienta, de 1 a 3 mm, y $\Delta$ la distancia adicional debida al ángulo de
posición $\omega$ de la herramienta: $\tan\omega = p/\Delta$. Si $\omega=90^\circ$
resulta $\Delta=0$ y si $\omega=45^\circ$ resulta $\Delta=p$.

\textbf{Rugosidad teórica.} Para una herramienta con radio de punta $r_\varepsilon$
y avance $a$, la rugosidad geométrica aproximada es~\refbib{bib:groover}
\begin{equation}\label{eq:Ra}
	R_a \approx \frac{a^2}{32\, r_\varepsilon}
\end{equation}

\textbf{Taladrado.} Para una broca de diámetro $d$ y ángulo de punta $\theta$, la
punta añade al recorrido una distancia
$\Delta = \dfrac{d/2}{\tan(\theta/2)}\approx 0{,}3\,d$ para $\theta=118^\circ$.

\subsection{Tolerancias y ajustes ISO}

El sistema ISO 286~\refbib{bib:iso286} expresa la tolerancia de una cota como
una letra (posición de la zona respecto de la línea de cero: mayúscula para
agujeros, minúscula para ejes) y un número (grado de tolerancia IT, que fija el
ancho de la zona). Para el agujero H la desviación inferior es $EI=0$ y para el
eje h la desviación superior es $es=0$. Los juegos extremos de un ajuste son
\begin{equation}\label{eq:juego}
	J_{\max}=ES-ei,\qquad J_{\min}=EI-es
\end{equation}
y si el resultado es negativo se trata de \emph{aprieto}. Los anchos de zona que
interesan en este trabajo se muestran en la Tabla~\ref{tab:IT}.

\begin{table}[!htb]
	\centering
	\caption{Anchos de tolerancia ISO utilizados (ISO 286-1)}\label{tab:IT}
	\small
	\rowcolors{2}{white}{lightblue}
	\begin{tabular}{c c c c}
		\rowcolor{darkblue}
		\encab{Rango nominal (mm)} & \encab{IT6 (\textmu m)} & \encab{IT7 (\textmu m)} & \encab{Cota del buje}    \\
		18 a 30                    & 13                      & 21                      & $\varnothing 25$ (H7/h6) \\
		30 a 50                    & 16                      & 25                      & $\varnothing 32$ (H7/p6) \\
	\end{tabular}
\end{table}

\subsection{Material: bronce al plomo SAE 660}

Los cojinetes de fricción se fabrican en bronces porque combinan resistencia
suficiente, buen comportamiento a la fricción y desgaste, y una dureza menor que
la del eje de acero, de modo que el desgaste se concentra en la pieza más barata
y fácil de reemplazar (el buje). El \textbf{SAE 660} (UNS C93200, ``bronce
83-7-7-3'') es el bronce al plomo-estaño más usado para casquillos de carga
media y velocidad moderada; el plomo forma inclusiones blandas que mejoran el
deslizamiento en caso de lubricación deficiente y que, además, fragmentan la
viruta, lo que lo hace muy fácil de mecanizar. La Tabla~\ref{tab:material}
resume sus propiedades típicas~\refbib{bib:astm}.

\begin{table}[!htb]
	\centering
	\caption{Propiedades típicas del bronce SAE 660 (UNS
		C93200)}\label{tab:material}
	\small
	\rowcolors{1}{white}{lightblue}
	\begin{tabularx}{\linewidth}{>{\bfseries}l X}
		Composición nominal (\%)    & Cu 83; Sn 7; Pb 7; Zn 3                                        \\
		Resistencia a la tracción   & $\approx 240$ MPa                                              \\
		Límite elástico             & $\approx 125$ MPa                                              \\
		Dureza                      & $\approx 65$--$70$ HB                                          \\
		Maquinabilidad              & $\approx 70\,\%$ (referencia: latón de corte libre $=100\,\%$) \\
		Densidad                    & $\approx 8{,}8$ g/cm$^3$                                       \\
		Coef. de dilatación térmica & $\approx 18\times10^{-6}\ ^\circ\text{C}^{-1}$                 \\
		Viruta                      & Corta y fragmentada (no requiere rompevirutas)                 \\
	\end{tabularx}
\end{table}

\subsection{Herramientas}

La Tabla~\ref{tab:herramientas} lista las herramientas del proceso. Para el
bronce (grupo ISO N, materiales no ferrosos) se emplea metal duro sin
recubrimiento y con geometría positiva y filo vivo, que corta con poca fuerza y
evita que el material se adhiera al filo. La broca se afila con un ángulo de
ataque reducido para que no se ``clave'' en el bronce.

\begin{table}[!htb]
	\centering
	\caption{Herramientas de la hoja de ruta}\label{tab:herramientas}
	\small
	\rowcolors{2}{white}{lightblue}
	\begin{tabularx}{\linewidth}{c l X}
		\rowcolor{darkblue}
		\encab{Cód.} & \encab{Herramienta}              & \encab{Descripción}                                                                                                       \\
		T01          & Cilindrado y refrentado exterior & Porta-inserto de $93^\circ$ con inserto rómbico $55^\circ$ tipo DCGT 11T304, metal duro ISO N, $r_\varepsilon=0{,}4$ mm   \\
		T02          & Broca de centros                 & HSS, $\varnothing 4$ mm                                                                                                   \\
		T03          & Broca helicoidal                 & HSS-Co, $\varnothing 22$ mm, punta $118^\circ$, con ataque reducido                                                       \\
		T04          & Mandrinado                       & Barra $\varnothing 16$ mm (mín. $\approx 20$ mm de agujero) con inserto tipo CCGT 060204, ISO N, $r_\varepsilon=0{,}4$ mm \\
		T05          & Tronzado                         & Porta-inserto de tronzar con inserto de 3 mm de ancho (tipo MGMN 300)                                                     \\
	\end{tabularx}
\end{table}

\newpage

% ===== 3. DESARROLLO =================================================
\section{DESARROLLO}

\subsection{Planos de la pieza}

La Figura~\ref{fig:plano} muestra el plano de fabricación del buje: corte
longitudinal con las cotas y tolerancias, y vista frontal. Las zonas rayadas son
la sección del material.

\begin{figure}[H]
	\centering
	\begin{tikzpicture}[scale=0.23,line cap=round,
			dim/.style={<->,>=Stealth,line width=0.4pt,color=black},
			ext/.style={line width=0.3pt,color=black!60},
			cl/.style={dash dot,line width=0.4pt,color=mediumblue},
			lead/.style={->,>=Stealth,line width=0.4pt,color=black},
			lbl/.style={font=\sffamily\scriptsize,text=darkblue}]
		% --- Sección: dos paredes simétricas respecto del eje
		\foreach \s in {1,-1}{
				\begin{scope}[yscale=\s]
					\path[pattern=north east lines,pattern color=mediumblue!70]
					(0,13)--(0,15)--(1,16)--(39,16)--(40,15)--(40,13)--(39.5,12.5)--(0.5,12.5)--cycle;
					\draw[draw=darkblue,line width=0.9pt,line join=round]
					(0,13)--(0,15)--(1,16)--(39,16)--(40,15)--(40,13)--(39.5,12.5)--(0.5,12.5)--cycle;
				\end{scope}}
		\draw[cl] (-4,0)--(44,0);
		% --- Longitud
		\draw[ext] (0,-17)--(0,-24); \draw[ext] (40,-17)--(40,-24);
		\draw[dim] (0,-22)--(40,-22) node[midway,below,lbl] {$40\pm0{,}1$};
		% --- Diámetro interior
		\draw[ext] (40.5,12.5)--(48,12.5); \draw[ext] (40.5,-12.5)--(48,-12.5);
		\draw[dim] (46,-12.5)--(46,12.5);
		\node[lbl,rotate=90,anchor=south] at (46,0) {$\varnothing25$ H7 $\left({}^{+0{,}021}_{\;\;0}\right)$};
		% --- Diámetro exterior
		\draw[ext] (41,16)--(55,16); \draw[ext] (41,-16)--(55,-16);
		\draw[dim] (53,-16)--(53,16);
		\node[lbl,rotate=90,anchor=south] at (53,0) {$\varnothing32$ p6 $\left({}^{+0{,}042}_{+0{,}026}\right)$};
		% --- Chaflanes
		\draw[lead] (-5,20)--(0.5,15.6); \node[lbl,anchor=east] at (-5,20) {$1\times45^\circ$};
		\draw[lead] (-5,7)--(0.25,12.8); \node[lbl,anchor=east] at (-5,7) {$0{,}5\times45^\circ$};
		% --- Rugosidad
		\draw[lead] (20,6)--(20,12.4); \node[lbl] at (20,4.5) {Ra 0,8 \textmu m};
		\draw[lead] (20,21)--(20,16.1); \node[lbl] at (20,22.4) {Ra 1,6 \textmu m};
		\node[lbl] at (20,-28) {\itshape Corte longitudinal (cotas en mm)};
	\end{tikzpicture}

	\vspace{6pt}
	\begin{minipage}[c]{0.34\linewidth}
		\centering
		\begin{tikzpicture}[scale=0.13,line cap=round,
				cl/.style={dash dot,line width=0.4pt,color=mediumblue},
				lbl/.style={font=\sffamily\scriptsize,text=darkblue}]
			\draw[draw=darkblue,line width=0.9pt,fill=lightblue] (0,0) circle (16);
			\draw[draw=darkblue,line width=0.9pt,fill=white] (0,0) circle (12.5);
			\draw[draw=darkblue,line width=0.4pt] (0,0) circle (15);
			\draw[draw=darkblue,line width=0.4pt] (0,0) circle (13);
			\draw[cl] (-19,0)--(19,0); \draw[cl] (0,-19)--(0,19);
			\node[lbl] at (0,-21.5) {\itshape Vista frontal};
		\end{tikzpicture}
	\end{minipage}\hfill
	\begin{minipage}[c]{0.62\linewidth}
		\small
		\rowcolors{1}{white}{lightblue}
		\begin{tabularx}{\linewidth}{>{\bfseries}l X}
			Pieza        & Buje de bronce                           \\
			Material     & SAE 660 (UNS C93200)                     \\
			Tol. general & ISO 2768-m                               \\
			Coaxialidad  & $\varnothing 0{,}02$ (agujero--exterior) \\
			Unidades     & mm; esquema no a escala                  \\
		\end{tabularx}
	\end{minipage}
	\caption{Plano de fabricación del buje: corte longitudinal y vista
		frontal.}\label{fig:plano}
\end{figure}

\begin{figure}[H]
	\centering
	\image[0.8][6cm]{img/plano_buje_cad.png}{Plano de fabricación en CAD (formato A4) con cajetín y vistas acotadas}
	\caption{Plano del buje elaborado en CAD (insertar el plano del grupo).}\label{fig:plano_cad}
\end{figure}

\subsection{Hoja de ruta}

El proceso se realiza en dos sujeciones de un torno CNC (Figura~\ref{fig:flujo}).
En la primera, la barra de $\varnothing 35$ mm queda en el plato con un voladizo
de 50 mm ($41$ mm de pieza $+\ 3$ mm del tronzador $+\ 6$ mm de holgura); se
refrenta, se hace el agujero y se tornea el exterior \emph{sin soltar la barra},
de modo que ambos diámetros quedan coaxiales por construcción. Después se
tronza la pieza a 41 mm y, en la segunda sujeción, se invierte en mordazas
blandas y se refrenta a la longitud final de 40 mm.

\begin{figure}[H]
	\centering
	\begin{tikzpicture}[node distance=5mm,
			caja/.style={rectangle,rounded
					corners=3pt,fill=darkblue,text=white,font=\sffamily\scriptsize\bfseries,
					minimum height=1.1cm,minimum width=2.3cm,align=center,inner sep=3pt},
			flecha/.style={-{Stealth[length=2.2mm]},thick,color=mediumblue}]
		\node[caja] (a) {Sujetar barra\\$\varnothing35$};
		\node[caja,right=of a] (b) {Refrentar\\cara 1};
		\node[caja,right=of b] (c) {Centrar y taladrar\\$\varnothing22$};
		\node[caja,right=of c] (d) {Mandrinar\\$\varnothing25$ H7};
		\node[caja,right=of d] (e) {Cilindrar\\$\varnothing32$ p6};
		\node[caja,below=7mm of a] (f) {Chaflanes\\cara 1};
		\node[caja,right=of f] (g) {Tronzar\\a 41 mm};
		\node[caja,right=of g] (h) {Invertir en\\mordazas blandas};
		\node[caja,right=of h] (i) {Refrentar a 40\\y chaflanes};
		\node[caja,fill=mediumblue,right=of i] (j) {Inspección\\final};
		\draw[flecha] (a)--(b); \draw[flecha] (b)--(c); \draw[flecha] (c)--(d); \draw[flecha] (d)--(e);
		\draw[flecha] (e.south) -- ++(0,-3.5mm) -| ($(f.north)+(0,3.5mm)$) -- (f.north);
		\draw[flecha] (f)--(g); \draw[flecha] (g)--(h); \draw[flecha] (h)--(i); \draw[flecha] (i)--(j);
	\end{tikzpicture}
	\caption{Secuencia general de fabricación del buje.}\label{fig:flujo}
\end{figure}

La Tabla~\ref{tab:ruta} es la hoja de ruta con las condiciones de corte
adoptadas (justificadas en la sección~\ref{sec:calculos}). Los tiempos de las
operaciones 20 a 110 son tiempos de corte; las operaciones 10, 90 y 120 son
tiempos auxiliares y se tratan aparte.

	{\small
		\setlength{\tabcolsep}{3pt}
		\rowcolors{2}{white}{lightblue}
		\begin{longtable}{c >{\raggedright\arraybackslash}p{3.8cm} c c c c c}
			\caption{Hoja de ruta del buje}\label{tab:ruta}
			\\
			\rowcolor{darkblue}
			\encab{Op.}                                        & \encab{Descripción}
			                                                   & \encab{Herr.}                                                           & \encab{$n$ (rpm)} &
			\encab{$a$ (mm/rev)}                               & \encab{$p$ (mm)}
			                                                   & \encab{$t$ (min)}
			\\
			\endfirsthead
			\rowcolor{darkblue}
			\encab{Op.}                                        & \encab{Descripción}                                                     & \encab{Herr.}     & \encab{$n$ (rpm)} &
			\encab{$a$ (mm/rev)}                               & \encab{$p$ (mm)}
			                                                   & \encab{$t$ (min)}
			\\
			\endhead
			10                                                 & Sujetar barra $\varnothing35$ en plato de 3 mordazas; voladizo de 50 mm & ---               & ---               & ---         & ---              & aux.  \\
			20                                                 & Refrentado cara 1: desbaste y acabado                                   & T01               & 2\,700 / 3\,400   & 0,30 / 0,10 & 0,7 / 0,3        & 0,081 \\
			30                                                 & Punto de centro                                                         & T02               & 2\,400            & 0,10        & ---              & 0,029 \\
			40                                                 & Taladrado $\varnothing22$, profundidad 46 mm                            & T03               & 430               & 0,20        & ---              & 0,635 \\
			50                                                 & Mandrinado a $\varnothing25{,}010$: 2 desbastes y 1 acabado             & T04               & 3\,000 / 4\,000   & 0,15 / 0,08 & 0,6 (×2) / 0,305 & 0,333 \\
			60                                                 & Cilindrado exterior a $\varnothing32{,}034$: 1 desbaste y 1 acabado     & T01               & 2\,700 / 3\,700   & 0,30 / 0,10 & 1,2 / 0,283      & 0,181 \\
			70                                                 & Chaflanes cara 1: ext. $1\times45^\circ$, int. $0{,}5\times45^\circ$    & T01, T04          & 3\,700 / 4\,000   & 0,10 / 0,08 & ---              & 0,012 \\
			80                                                 & Tronzado a 41 mm de
			la cara 1                                          & T05                                                                     & 2\,700
			                                                   & 0,08                                                                    & ---               & 0,026                                                      \\
			90                                                 & Invertir la pieza en mordazas blandas                                   & ---               & ---               & ---         & ---              & aux.  \\
			100                                                & Refrentado cara 2 a 40 mm: desbaste y acabado                           & T01               & 2\,700 / 3\,700   & 0,30 / 0,10 & 0,7 / 0,3        & 0,022 \\
			110                                                & Chaflanes cara 2 (iguales a la op. 70)                                  & T01, T04
			                                                   & 3\,700 / 4\,000                                                         & 0,10 / 0,08       & ---               & 0,012                                  \\
			120                                                & Desbarbado e inspección final                                           & ---               & ---               & ---         & ---              & aux.  \\
			\hline
			\multicolumn{6}{r}{\textbf{Tiempo de corte total}} & \textbf{1,331}
			\\
		\end{longtable}}

\subsection{Cálculos de parámetros}\label{sec:calculos}

\subsubsection{Selección de las velocidades de corte}

Para el bronce con inserto de metal duro la tabla del curso (Anexo~A) recomienda
$v=300$ m/min en desbaste y $v=380$ m/min en acabado; para herramientas de acero
rápido (broca de centros y broca helicoidal) recomienda $v=30$ m/min. El
procedimiento en todas las operaciones es el mismo: (1) obtener $n$ ideal con
$n=1000\,v/(\pi D_i)$; (2) adoptar una $n$ realizable por la máquina, sin superar
la $v$ de la tabla; (3) recalcular la $v$ real con la ecuación~\eqref{eq:v}.

\subsubsection{Cilindrado exterior (OP 60)}

Se parte de $D_i=35$ mm. El diámetro final se fija en el \textbf{centro} de la
zona p6 ($32{,}026$ a $32{,}042$ mm), es decir $D_f=32{,}034$ mm, para repartir
el margen de error hacia ambos lados.
\begin{align*}
	p_T & = \frac{D_i-D_f}{2}=\frac{35-32{,}034}{2}=1{,}483\umm                                   \\
	p_T & = N_DP_D+N_AP_A\ \Rightarrow\ 1{,}483 = 1(1{,}2)+1\,(P_A)\ \Rightarrow\ P_A=0{,}283\umm
\end{align*}
Se adopta una pasada de desbaste de $1{,}2$ mm y una de acabado de $0{,}283$ mm
(dentro del intervalo 0,1--0,5 mm). El diámetro inicial del acabado es
$D_{i,A}=35-2(1{,}2)=32{,}6$ mm.

\textbf{Velocidad de giro y de corte.}
\begin{align*}
	n_D^{\text{ideal}} & =\frac{1000(300)}{\pi(35)}=2\,728\urpm\ \Rightarrow\ n_D=2\,700\urpm     \\
	v_D                & =\frac{\pi(35)(2\,700)}{1000}=296{,}9\umpm                               \\
	n_A^{\text{ideal}} & =\frac{1000(380)}{\pi(32{,}6)}=3\,710\urpm\ \Rightarrow\ n_A=3\,700\urpm \\
	v_A                & =\frac{\pi(32{,}6)(3\,700)}{1000}=378{,}9\umpm
\end{align*}
Ambas velocidades quedan por debajo de las de la tabla. Se cumple la regla
$v_D<v_A$, $a_D>a_A$, $n_D<n_A$.

\textbf{Avance y velocidad lineal.} Con $a_D=0{,}30$ mm/rev y $a_A=0{,}10$ mm/rev:
\begin{align*}
	v_{L,D} & =a_Dn_D=0{,}30(2\,700)=810\ummmin \\
	v_{L,A} & =a_An_A=0{,}10(3\,700)=370\ummmin
\end{align*}

\textbf{Tiempo de pasada.} Se mecaniza $L=44$ mm (41 mm de pieza más 3 mm del
ancho del tronzador), con $y=2$ mm y, al ser $\omega=90^\circ$, $\Delta=0$:
\[
	t_D=\frac{44+2+0}{810}=0{,}0568\umin\ (3{,}4\ \text{s}),\qquad
	t_A=\frac{44+2+0}{370}=0{,}1243\umin\ (7{,}5\ \text{s})
\]

\textbf{Caudal de viruta.}
\begin{align*}
	A_D & =\tfrac{\pi}{4}\left(35^2-32{,}6^2\right)=127{,}4\ \text{mm}^2      \\
	Q_D & =810(127{,}4)=103\,213\ \text{mm}^3/\text{min}=103{,}2\ucmmin       \\
	A_A & =\tfrac{\pi}{4}\left(32{,}6^2-32{,}034^2\right)=28{,}7\ \text{mm}^2 \\
	Q_A & =370(28{,}7)=10\,632\ \text{mm}^3/\text{min}=10{,}6\ucmmin
\end{align*}

\textbf{Rugosidad teórica del acabado} (ec.~\ref{eq:Ra}, $r_\varepsilon=0{,}4$
mm):
$R_a\approx 0{,}10^2/(32\cdot0{,}4)=0{,}00078$ mm $=0{,}78$ \textmu m, inferior
al
límite de 1,6 \textmu m.

\subsubsection{Mandrinado o cilindrado interior (OP 50)}

El agujero parte de la broca ($D_i=22$ mm) y debe llegar al centro de la zona
H7 ($25{,}000$ a $25{,}021$ mm): $D_f=25{,}010$ mm.
\begin{align*}
	p_T & =\frac{25{,}010-22}{2}=1{,}505\umm                                                   \\
	p_T & =N_DP_D+N_AP_A\ \Rightarrow\ 1{,}505=2(0{,}6)+1\,(P_A)\ \Rightarrow\ P_A=0{,}305\umm
\end{align*}
Los diámetros resultantes son $22\to23{,}2\to24{,}4\to25{,}010$ mm, con
$D_{i,A}=24{,}4$ mm. Con la tabla, las rpm ideales serían 4\,341 rpm
(desbaste) y 4\,957 rpm (acabado), que superan el máximo de la máquina
(4\,000 rpm) y exigirían una barra más rígida. Se adopta $n_D=3\,000$ rpm, para
dejar margen de rigidez a la barra y de evacuación de viruta, y $n_A=4\,000$ rpm:
\[
	v_D=\frac{\pi(22)(3\,000)}{1000}=207{,}3\umpm,\qquad
	v_A=\frac{\pi(24{,}4)(4\,000)}{1000}=306{,}6\umpm
\]
Con $a_D=0{,}15$ mm/rev y $a_A=0{,}08$ mm/rev (menores que en el exterior por
el voladizo de la barra): $v_{L,D}=450$ mm/min y $v_{L,A}=320$ mm/min. Con $L=42$
mm (se mandrina hasta pasar el plano de tronzado), $y=2$ mm y $\Delta=0$:
\[
	t_D=N_D\left(\frac{42+2}{450}\right)=2(0{,}0978)=0{,}1956\umin,\qquad
	t_A=\frac{42+2}{320}=0{,}1375\umin
\]
\textbf{Caudal de viruta:} primera pasada de desbaste
$A=\frac{\pi}{4}(23{,}2^2-22^2)=42{,}6$ mm$^2$ y $Q=450(42{,}6)=19{,}2\ucmmin$;
segunda pasada $A=44{,}9$ mm$^2$ y $Q=20{,}2\ucmmin$; acabado
$A=\frac{\pi}{4}(25{,}010^2-24{,}4^2)=23{,}7$ mm$^2$ y
$Q=320(23{,}7)=7{,}6\ucmmin$. Rugosidad teórica del acabado:
$R_a\approx0{,}08^2/(32\cdot0{,}4)=0{,}5$ \textmu m $<0{,}8$ \textmu m.

\subsubsection{Punto de centro y taladrado (OP 30 y 40)}

Ambas herramientas son de acero rápido, por lo que se usa $v=30$ m/min.
\begin{align*}
	n      & =\frac{1000(30)}{\pi(22)}=434\urpm\ \Rightarrow\ n=430\urpm                      \\
	v      & =\frac{\pi(22)(430)}{1000}=29{,}7\umpm                                           \\
	\Delta & =\frac{22/2}{\tan 59^\circ}=6{,}61\umm,\qquad v_L=0{,}20(430)=86\ummmin          \\
	t      & =\frac{L+y+\Delta}{a\,n}=\frac{46+2+6{,}61}{86}=0{,}635\umin\ (38{,}1\ \text{s})
\end{align*}
La profundidad de 46 mm asegura que el agujero pase el plano de tronzado (a 41 a
44 mm de la cara). El caudal de viruta del taladrado es
$Q=v_L\,\frac{\pi}{4}d^2=86(380{,}1)=32{,}7\ucmmin$. Para el punto de centro
($\varnothing4$ mm): $n=1000(30)/(\pi\cdot4)=2\,387\to2\,400$ rpm, $a=0{,}10$
mm/rev, $L=5$ mm, $y=2$ mm: $t=7/(0{,}10\cdot2\,400)=0{,}029$ min.

\subsubsection{Refrentados (OP 20 y 100)}

En cada cara se quitan $1{,}0$ mm axiales: $0{,}7$ mm en desbaste y $0{,}3$ mm en
acabado. En el refrentado de la cara 1 (barra maciza) la herramienta recorre
$L=D/2=17{,}5$ mm; para que $v\le380$ m/min en el acabado se adopta
$n_A=3\,400$ rpm ($v=373{,}8$ m/min).
\[
	t_D=\frac{17{,}5+2+0}{0{,}30(2\,700)}=0{,}0241\umin,\qquad
	t_A=\frac{17{,}5+2+0}{0{,}10(3\,400)}=0{,}0574\umin
\]
En la cara 2 la pieza ya es hueca y la herramienta solo recorre la pared:
$L=(32{,}034-25{,}010)/2=3{,}512$ mm, con $n_D=2\,700$ y $n_A=3\,700$ rpm:
\[
	t_D=\frac{3{,}512+2}{810}=0{,}0068\umin,\qquad
	t_A=\frac{3{,}512+2}{370}=0{,}0149\umin
\]
La longitud final queda $41-0{,}7-0{,}3=40{,}0$ mm, centrada en $40\pm0{,}1$.

\subsubsection{Tronzado (OP 80)}

Se corta una pared de $L=3{,}512$ mm (pieza hueca), con $y=2$ mm, $\Delta=0$,
$a=0{,}08$ mm/rev y $n=2\,700$ rpm ($v=271{,}7$ m/min en el diámetro exterior):
\[
	t=\frac{3{,}512+2+0}{0{,}08(2\,700)}=\frac{5{,}512}{216}=0{,}0255\umin\ (1{,}5\ \text{s})
\]

\subsubsection{Chaflanes (OP 70 y 110)}

Cada chaflán es una pasada recta a $45^\circ$ de longitud $\ell=\sqrt2\,c$. Con
$y=1$ mm: exterior $c=1$ mm, $a=0{,}10$ mm/rev, $n=3\,700$ rpm,
$t=(1{,}414+1)/370=0{,}0065$ min; interior $c=0{,}5$ mm, $a=0{,}08$ mm/rev,
$n=4\,000$ rpm, $t=(0{,}707+1)/320=0{,}0053$ min. Por cara resultan $0{,}012$
min
y entre ambas caras $0{,}024$ min.

\subsubsection{Resumen de condiciones y tiempos}

La Tabla~\ref{tab:resumen} reúne los resultados de las operaciones de corte.

\begin{table}[!htb]
	\centering
	\caption{Resumen de condiciones de corte, caudales y tiempos}\label{tab:resumen}
	\footnotesize
	\setlength{\tabcolsep}{3pt}
	\rowcolors{2}{white}{lightblue}
	\begin{tabular}{l c c c c r}
		\rowcolor{darkblue}
		\encab{Pasada}                   & \encab{$v$ (m/min)} & \encab{$n$ (rpm)} & \encab{$a$ (mm/rev)} & \encab{$Q$ (cm$^3$/min)} & \encab{$t$ (min)} \\
		Refrentado cara 1, desbaste      & 296,9               & 2\,700            & 0,30                 & ---                      & 0,024             \\
		Refrentado cara 1, acabado       & 373,8               & 3\,400            & 0,10                 & ---                      & 0,057             \\
		Punto de centro                  & 30,2                & 2\,400            & 0,10                 & ---                      & 0,029             \\
		Taladrado $\varnothing22$        & 29,7                & 430               & 0,20                 & 32,7                     & 0,635             \\
		Mandrinado, desbaste (2 pasadas) & 207,3               & 3\,000            & 0,15                 & 19,2 / 20,2              & 0,196             \\
		Mandrinado, acabado              & 306,6               & 4\,000            & 0,08                 & 7,6                      & 0,138             \\
		Cilindrado exterior, desbaste    & 296,9               & 2\,700            & 0,30                 & 103,2                    & 0,057             \\
		Cilindrado exterior, acabado     & 378,9               & 3\,700            & 0,10                 & 10,6                     & 0,124             \\
		Chaflanes (cara 1)               & ---                 & ---               &
		---                              & ---                 & 0,012                                                                                   \\
		Tronzado                         & 271,7               & 2\,700            & 0,08                 & ---                      & 0,026             \\
		Refrentado cara 2, desbaste      & 271,7               & 2\,700            & 0,30                 & ---                      & 0,007             \\
		Refrentado cara 2, acabado       & 372,4               & 3\,700            & 0,10                 & ---                      & 0,015             \\
		Chaflanes (cara 2)               & ---                 & ---               & ---                  & ---                      & 0,012             \\
		\midrule
		\textbf{Tiempo de corte total}   &                     &                   &                      &                          &
		\textbf{1,331}                                                                                                                                   \\
	\end{tabular}
\end{table}

\textbf{Tiempo de ciclo estimado.} A los $1{,}331$ min de corte se suman los
tiempos auxiliares estimados: carga y descarga (2 sujeciones $\times$ 0,75 min
$=1{,}5$ min), 7 cambios de herramienta ($\times$ 0,10 min $=0{,}7$ min) e
inspección (1,0 min). El ciclo es
\[
	t_{\text{ciclo}}=1{,}331+1{,}5+0{,}7+1{,}0\approx4{,}53\umin
	\ \Rightarrow\ \frac{60}{4{,}53}\approx13\ \text{piezas/h}
\]

\subsection{Registro gráfico del proceso}

Al ser un trabajo teórico, las imágenes de la Figura~\ref{fig:fotos} son
referenciales (se reemplazarán por fotos o ilustraciones del grupo, indicando la
fuente).

\begin{figure}[H]
	\centering
	\begin{minipage}[t]{0.48\linewidth}\centering
		\image[0.97][3.6cm]{img/proc_refrentado.jpg}{Refrentado de la cara 1 de la barra}\\[2pt]
		{\sffamily\scriptsize (a) Refrentado}
	\end{minipage}\hfill
	\begin{minipage}[t]{0.48\linewidth}\centering
		\image[0.97][3.6cm]{img/proc_taladrado.jpg}{Taladrado con broca $\varnothing22$}\\[2pt]
		{\sffamily\scriptsize (b) Taladrado}
	\end{minipage}\\[8pt]
	\begin{minipage}[t]{0.48\linewidth}\centering
		\image[0.97][3.6cm]{img/proc_mandrinado.jpg}{Mandrinado con barra de inserto}\\[2pt]
		{\sffamily\scriptsize (c) Mandrinado}
	\end{minipage}\hfill
	\begin{minipage}[t]{0.48\linewidth}\centering
		\image[0.97][3.6cm]{img/proc_cilindrado.jpg}{Cilindrado exterior de acabado}\\[2pt]
		{\sffamily\scriptsize (d) Cilindrado exterior}
	\end{minipage}\\[8pt]
	\begin{minipage}[t]{0.48\linewidth}\centering
		\image[0.97][3.6cm]{img/proc_tronzado.jpg}{Tronzado de la pieza}\\[2pt]
		{\sffamily\scriptsize (e) Tronzado}
	\end{minipage}\hfill
	\begin{minipage}[t]{0.48\linewidth}\centering
		\image[0.97][3.6cm]{img/buje_terminado.jpg}{Buje terminado y tapón pasa/no pasa}\\[2pt]
		{\sffamily\scriptsize (f) Buje terminado}
	\end{minipage}
	\caption{Secuencia de operaciones del buje (imágenes
		referenciales).}\label{fig:fotos}
\end{figure}

\newpage

% ===== 4. RESULTADOS =================================================
\section{RESULTADOS}

Como el trabajo es teórico, los resultados son los valores que el proceso
planificado debe producir y la comprobación analítica de que estos cumplen el
plano. No se incluyen lecturas de instrumentos reales; la tabla de inspección
está diseñada para que el grupo pueda añadir una columna de ``valor medido''
cuando se fabrique la pieza.

\subsection{Tabla de mediciones}

La Tabla~\ref{tab:pasadas} recoge el diámetro teórico que debe medirse después
de
cada pasada, calculado con las profundidades de la sección~\ref{sec:calculos}.
En el agujero se mide después de cada pasada de desbaste solo si se desea
controlar el proceso; la medición obligatoria es la del acabado.

\begin{table}[!htb]
	\centering
	\caption{Dimensiones teóricas esperadas después de cada pasada}\label{tab:pasadas}
	\small
	\rowcolors{2}{white}{lightblue}
	\begin{tabular}{l l c c}
		\rowcolor{darkblue}
		\encab{Característica} & \encab{Pasada}              & \encab{Prof. radial (mm)} & \encab{Dimensión (mm)} \\
		Agujero                & Broca                       & ---                       & 22,000                 \\
		Agujero                & Desbaste 1                  & 0,6                       & 23,200                 \\
		Agujero                & Desbaste 2                  & 0,6                       & 24,400                 \\
		Agujero                & \textbf{Acabado}            & 0,305
		                       & \textbf{25,010}                                                                  \\
		Exterior               & Barra                       & ---
		                       & 35,000                                                                           \\
		Exterior               & Desbaste                    & 1,2
		                       & 32,600                                                                           \\
		Exterior               & \textbf{Acabado}            & 0,283
		                       & \textbf{32,034}
		\\
		Longitud               & Tronzado                    & ---
		                       & 41,000                                                                           \\
		Longitud               & Refrentado desbaste         & 0,7 (axial)
		                       & 40,300                                                                           \\
		Longitud               & \textbf{Refrentado acabado} & 0,3 (axial)
		                       & \textbf{40,000}                                                                  \\
	\end{tabular}
\end{table}

La Tabla~\ref{tab:inspeccion} es el plan de inspección. La columna final es la
relación entre el ancho de la tolerancia y la resolución del instrumento: el
criterio habitual en metrología dimensional es que el instrumento resuelva al
menos una décima parte de la tolerancia (relación $\ge10$).

\begin{table}[!htb]
	\centering
	\caption{Plan de inspección y límites de aceptación}\label{tab:inspeccion}
	\footnotesize
	\setlength{\tabcolsep}{3pt}
	\rowcolors{2}{white}{lightblue}
	\begin{tabularx}{\linewidth}{>{\raggedright\arraybackslash}p{2.1cm} c c >{\raggedright\arraybackslash}X c c}
		\rowcolor{darkblue}
		\encab{Característica}                                      & \encab{Límites (mm)}              & \encab{Objetivo (mm)} & \encab{Instrumento}                            & \encab{Resol. (mm)} & \encab{Tol./resol.} \\
		$\varnothing25$ H7                                          & 25,000 -- 25,021                  & 25,010                &
		Micrómetro de interiores de 3 puntos; tapón pasa/no pasa H7 & 0,001                             & 21                                                                                                                 \\
		$\varnothing32$ p6                                          & 32,026 -- 32,042                  & 32,034                & Micrómetro de exteriores 25--50 mm             & 0,001               & 16                  \\
		Longitud 40                                                 & 39,9 -- 40,1                      & 40,0                  & Vernier                                        & 0,02                & 10                  \\
		Coaxialidad                                                 & $\le0{,}02$                       & ---                   & Comparador de carátula sobre mandril expansivo & 0,001               & 20                  \\
		Ra agujero / exterior                                       & $\le0{,}8$ / $\le1{,}6$ \textmu m & ---                   & Rugosímetro                                    & ---                 & ---                 \\
	\end{tabularx}
\end{table}

Aquí se aprecia la conexión con el laboratorio de metrología: un vernier de
0,05 mm de resolución habría dado solo una relación de 0,42 sobre la zona H7
(21 \textmu m), por lo que no sirve para verificar este agujero. Solo el
micrómetro de interiores (o el calibre pasa/no pasa) es adecuado.

\subsection{Verificación de tolerancias}

\textbf{Ajuste H7/h6 entre el agujero y el eje.} Con $\varnothing25$ H7
($ES=+21$ \textmu m, $EI=0$) y h6 ($es=0$, $ei=-13$ \textmu m):
\begin{align*}
	D_{\max} & =25{,}021\umm,          & D_{\min} & =25{,}000\umm \\
	d_{\max} & =25{,}000\umm,          & d_{\min} & =24{,}987\umm \\
	J_{\max} & =ES-ei=21-(-13)=34\uum, & J_{\min} & =EI-es=0-0=0
\end{align*}
Es un ajuste con juego, que va de 0 a 34 \textmu m (Figura~\ref{fig:zonas}). Si el
agujero se mecaniza en el valor objetivo de $25{,}010$ mm, el juego queda entre
$25{,}010-25{,}000=10$ \textmu m y $25{,}010-24{,}987=23$ \textmu m, es decir,
siempre positivo, lo que protege contra el agarrotamiento aunque el eje salga
en cualquier punto de su zona.

\begin{figure}[H]
	\centering
	\begin{tikzpicture}[x=0.9cm,y=0.115cm,
			lbl/.style={font=\sffamily\scriptsize,text=darkblue}]
		\draw[->,line width=0.5pt] (-0.4,-18)--(-0.4,26) node[above,lbl] {\textmu m};
		\foreach \yy in {-13,0,10,21}{
				\draw (-0.5,\yy)--(-0.3,\yy); \node[lbl,anchor=east] at (-0.5,\yy) {\yy};}
		\draw[dash dot,line width=0.5pt,color=black!70] (-0.4,0)--(10.8,0)
		node[right,lbl,align=left] {línea de cero\\$\varnothing25{,}000$};
		\filldraw[fill=lightblue,draw=darkblue,line width=0.8pt] (0.8,0) rectangle (3.8,21);
		\node[lbl,align=center] at (2.3,29.5) {Agujero H7\\[-1pt]$ES=+21$, $EI=0$};
		\draw[dashed,color=mediumblue,line width=0.6pt] (0.8,10)--(3.8,10);
		\node[lbl,fill=lightblue,inner sep=1pt] at (2.3,14) {objetivo CNC $+10$};
		\filldraw[fill=lightblue,draw=darkblue,line width=0.8pt] (5.6,-13) rectangle (8.6,0);
		\node[lbl,align=center] at (7.1,-23) {Eje h6\\[-1pt]$es=0$, $ei=-13$};
		\draw[<->,>=Stealth,line width=0.5pt] (9.6,-13)--(9.6,21);
		\node[lbl,rotate=90,anchor=south,xshift=-2pt] at (9.6,4) {$J_{\max}=34\
				\mu$m};
	\end{tikzpicture}
	\caption{Zonas de tolerancia del ajuste $\varnothing25$ H7/h6 (valores en
		\textmu m respecto de la medida nominal).}\label{fig:zonas}
\end{figure}

\textbf{Ajuste exterior H7/p6.} Con el alojamiento $\varnothing32$ H7
($32{,}000$ a $32{,}025$ mm) y el buje p6 ($32{,}026$ a $32{,}042$ mm) el
aprieto
va de $32{,}026-32{,}025=1$ \textmu m a $32{,}042-32{,}000=42$ \textmu m. Con el
objetivo de $32{,}034$ mm queda entre 9 y 34 \textmu m, siempre positivo, de
modo
que el buje no gira en la carcasa.

\textbf{Longitud, rugosidad y geometría.} La longitud nominal resulta
$41-0{,}7-0{,}3=40{,}0$ mm, centrada en la banda de $\pm0{,}1$ mm. Las
rugosidades teóricas de acabado ($0{,}5$ \textmu m en el agujero y $0{,}78$ \textmu m
en el exterior) son menores que los límites de $0{,}8$ y $1{,}6$ \textmu m. La
coaxialidad se obtiene al mecanizar agujero y exterior sin soltar la barra, de
modo que depende solo de la precisión del husillo y del conjunto, no de un
reposicionado.

\subsection{Análisis de errores}

El diámetro del agujero H7 admite una dispersión total de 21 \textmu m, es decir,
$\pm10{,}5$ \textmu m alrededor del objetivo. La Tabla~\ref{tab:errores} estima las
principales fuentes de error \emph{diametral} del mandrinado de acabado.

\begin{table}[!htb]
	\centering
	\caption{Presupuesto de error estimado del agujero $\varnothing25$ H7}\label{tab:errores}
	\small
	\rowcolors{2}{white}{lightblue}
	\begin{tabularx}{\linewidth}{>{\raggedright\arraybackslash}p{4.0cm} X r}
		\rowcolor{darkblue}
		\encab{Fuente}                       & \encab{Estimación}                                                                                                                                                                                      & \encab{Error diam. (\textmu m)} \\
		Dilatación térmica                   & $\alpha D\Delta T=18\times10^{-6}(25)(5)$, con $\Delta T=5$ K entre mecanizado y medición                                                                                                               & 2,25                            \\
		Deflexión de la barra                & $F_c=k_sa\,p=700(0{,}08)(0{,}305)\approx17$ N; $F_r\approx0{,}4F_c\approx7$ N; $\delta=\dfrac{F_rL^3}{3EI}=0{,}73$ \textmu m con barra de acero $\varnothing16$, voladizo de 60 mm; diámetro $=2\delta$ & 1,5                             \\
		Desgaste de la herramienta           & 3 \textmu m de desgaste radial entre correcciones (supuesto); afecta dos veces al diámetro                                                                                                              & 6,0                             \\
		Incertidumbre de medición            & Micrómetro de interiores de 3 puntos (supuesto)                                                                                                                                                         & 3,0                             \\
		\midrule
		\textbf{Suma cuadrática (RSS)}       & $\sqrt{2{,}25^2+1{,}5^2+6^2+3^2}$                                                                                                                                                                       & \textbf{7,2}                    \\
		\textbf{Suma aritmética (peor caso)} & $2{,}25+1{,}5+6+3$                                                                                                                                                                                      & \textbf{12,75}                  \\
	\end{tabularx}
\end{table}

Lo que sale de la tabla es esto:
\begin{itemize}
	\item \textbf{El desgaste de la herramienta domina el presupuesto.} Un
	      desgaste radial de apenas 3 \textmu m cambia el diámetro 6 \textmu m,
	      casi un tercio de toda la zona H7. Es la razón de fondo por la que el
	      objetivo se fija en el centro de la zona y no cerca de su límite
	      inferior.
	\item \textbf{El error típico cabe en la tolerancia, el peor caso no.} La
	      combinación cuadrática (7,2 \textmu m) deja un margen de 3,3 \textmu m
	      respecto de la semibanda de 10,5 \textmu m; la suma aritmética
	      (12,75 \textmu m) la supera. Por tanto, el proceso es viable pero
	      \emph{no} es robusto sin control: hace falta medir y corregir.
	\item \textbf{El exterior p6 tiene menos margen.} Con $\alpha D\Delta T=2{,}9$
	      \textmu m, desgaste de 6 \textmu m y medición de 3 \textmu m, la suma
	      cuadrática es 7,3 \textmu m frente a una semibanda de 8 \textmu m
	      (IT6 = 16 \textmu m): el margen es de solo 0,7 \textmu m.
	\item \textbf{Otras fuentes no cuantificadas:} deformación del buje por la
	      presión de las mordazas en la sujeción 2 (no afecta al agujero, que ya
	      está terminado, pero sí a la planitud de la cara 2); vibración de la
	      barra de mandrinar; viruta atrapada en el agujero; y rebabas del
	      tronzado.
\end{itemize}

\textbf{Estrategia de control propuesta.} Mecanizar una primera pieza y medirla
con micrómetro de interiores; corregir el desplazamiento de herramienta
(\emph{offset}) con la diferencia entre el valor medido y el objetivo de
25,010 mm; medir una pieza cada cierto número de piezas y corregir el desgaste
antes de que el diámetro se acerque a los límites; y verificar con el tapón
pasa/no pasa H7 en el 100\,\% de las piezas.

\subsection{Productividad y distribución del tiempo}

El tiempo de corte total es de 1,331 min (79,8 s) y el ciclo estimado, de unos
4,53 min por pieza. La Figura~\ref{fig:tiempos} muestra que el taladrado
consume casi la mitad del tiempo de corte (47,7\,\%) y el mandrinado otro cuarto
(25,0\,\%): entre los dos suman el 72,7\,\%. Las operaciones de torneado exterior,
refrentado, tronzado y chaflanes, que son las del curso, ocupan una fracción
pequeña del ciclo.

\begin{figure}[H]
	\centering
	\begin{tikzpicture}
		\begin{axis}[xbar,/pgf/number format/use comma,width=0.8\linewidth,height=6.6cm,
				symbolic y coords={A,B,C,D,E,F,G,H},ytick=data,y dir=reverse,
				yticklabels={Taladrado,Mandrinado,Cilindrado ext.,Refrentado cara 1,Punto de centro,Tronzado,Chaflanes,Refrentado cara 2},
				xmin=0,xmax=0.82,bar width=7pt,enlarge y limits=0.08,
				xlabel={Tiempo de corte (min)},
				xmajorgrids,grid style={borderblue!70},
				label style={font=\sffamily\small},tick label style={font=\sffamily\small},
				nodes near coords,
				nodes near coords style={font=\sffamily\scriptsize,/pgf/number format/.cd,use comma,fixed,precision=3,fixed zerofill},
				axis x line*=bottom,axis y line*=left]
			\addplot[fill=mediumblue,draw=darkblue] coordinates {(0.635,A) (0.333,B)
					(0.181,C) (0.081,D) (0.029,E) (0.026,F) (0.024,G) (0.022,H)};
		\end{axis}
	\end{tikzpicture}
	\caption{Distribución del tiempo de corte por operación (total 1,331 min).}\label{fig:tiempos}
\end{figure}

El material de partida también pesa: por pieza se consume $1+41+3=45$ mm de
barra
$\varnothing35$ ($V\approx43{,}3$ cm$^3$, $\approx0{,}38$ kg), de los cuales
queda en la pieza $12{,}5$ cm$^3$: un aprovechamiento del $29\,\%$. El resto es
viruta, en su mayor parte el núcleo del agujero.

% ===== 5. CONCLUSIONES ===============================================
\newpage
\section{CONCLUSIONES}

\subsection{Logros y dificultades}

\begin{itemize}
	\item Se definió una secuencia completa de fabricación del buje en dos
	      sujeciones y once operaciones, con las fórmulas del curso aplicadas a
	      cada una (velocidad de corte, rpm, avance, número de pasadas, caudal de
	      viruta y tiempos). El tiempo de corte resultó de 1,331 min y el ciclo
	      estimado de 4,53 min por pieza.\\[-10pt]
	\item El ajuste H7/h6 queda garantizado en el papel: juego entre 0 y 34 \textmu m,
	      que con el objetivo de $25{,}010$ mm se reduce a 10--23 \textmu m, siempre
	      positivo. Los instrumentos propuestos cumplen la relación
	      tolerancia/resolución de al menos 10:1.\\[-10pt]
	\item Hacer agujero y exterior en la misma sujeción de la barra asegura la
	      coaxialidad sin depender de un reposicionado, lo que justifica el uso de
	      un torno CNC frente a uno paralelo.\\[-10pt]
	\item La principal dificultad es de \emph{control}, no de cálculo: el
	      presupuesto de error muestra que el proceso cumple H7 en el caso típico
	      (7,2 \textmu m frente a 10,5 \textmu m), pero no en el peor caso
	      (12,75 \textmu m), y que el desgaste de la herramienta es la fuente
	      dominante. El exterior p6 tiene un margen aún menor.\\[-10pt]
	\item La tabla de velocidades del curso (300 y 380 m/min en bronce con
	      carburo) exige 4\,300 a 5\,000 rpm en el agujero de $\varnothing22$ a
	      $\varnothing25$ mm, por encima del límite supuesto de la máquina. Hubo que
	      reducir las rpm y apoyarse en el criterio de rigidez de la barra.\\[-10pt]
	\item Limitación del trabajo: todos los valores son teóricos; la
	      maquinabilidad, el desgaste, las rigideces y los tiempos auxiliares son
	      supuestos que solo una prueba real puede confirmar.
\end{itemize}

\subsection{Mejoras propuestas}

\begin{itemize}
	\item \textbf{Usar barra hueca (tubo de bronce) en lugar de barra maciza.} Se
	      elimina el taladrado ($48\,\%$ del tiempo de corte) y se aprovecha mucho
	      más material; es la mejora de mayor impacto.\\[-10pt]
	\item \textbf{Escariar el agujero con un escariador H7} en lugar de acabarlo
	      con barra: el diámetro lo fija directamente la herramienta y no depende
	      de compensar un filo que se desgasta, que es la mayor contribución del
	      presupuesto de errores.\\[-10pt]
	\item \textbf{Taladrar con ciclo de picado y refrigerante} para evacuar la
	      viruta en un agujero de profundidad $2{,}1\,d$, y revisar el avance de la
	      broca si el tiempo de ciclo es crítico.\\[-10pt]
	\item \textbf{Velocidad de corte constante (G96) en los refrentados}, para que
	      la velocidad de corte no caiga hacia el centro; los cálculos de este
	      trabajo usan rpm fijas, que son conservadoras en tiempo.\\[-10pt]
	\item \textbf{Compensación automática del desgaste} mediante palpador de
	      herramienta o medición en máquina, para reducir la contribución del
	      desgaste al error.\\[-10pt]
	\item \textbf{Prever el cierre del agujero por el aprieto del montaje.} Al
	      prensar un buje con aprieto en un alojamiento, su diámetro interior
	      suele
	      reducirse; si el ajuste H7/h6 es crítico se deja una sobremedida de
	      acabado y se mandrina o escaría el agujero \emph{después} del montaje.
	      \\[-10pt]
	\item \textbf{Fabricar una pieza de prueba} y reemplazar los valores
	      supuestos de desgaste, tiempos auxiliares y rigidez por valores medidos.
\end{itemize}

% ===== 6. ANEXOS =====================================================
\newpage
\section{ANEXOS}

\subsection*{Anexo A. Tabla de velocidades de corte ($v$, m/min)}

Fuente: material de clase de la Semana~4~\refbib{bib:clase}. Las filas utilizadas
en este trabajo son las del bronce.

\begin{table}[!htb]
	\centering
	\small
	\rowcolors{3}{white}{lightblue}
	\begin{tabular}{l c c c c}
		\rowcolor{darkblue}
		                             & \multicolumn{2}{c}{\encab{Acero rápido}} &
		\multicolumn{2}{c}{\encab{Pastilla carburada}}
		\\
		\rowcolor{darkblue}
		\encab{Material de la pieza} & \encab{Desbaste}                         &
		\encab{Acabado}              &
		\encab{Desbaste}             & \encab{Acabado}
		\\
		Acero 0,35 \%C               & 25                                       & 30
		                             & 200                                      & 300                              \\
		Acero 0,45 \%C               & 15                                       & 20
		                             & 120                                      & 160                              \\
		Acero extra duro             & 12                                       & 16
		                             & 40                                       & 60                               \\
		Hierro fundido maleable      & 20                                       & 25
		                             & 70                                       & 85                               \\
		Hierro fundido gris          & 15                                       & 20
		                             & 65                                       & 95                               \\
		Hierro fundido duro          & 10                                       & 15          & 30           & 50  \\
		\textbf{Bronce}              & \textbf{30}                              & \textbf{40} & \textbf{300} &
		\textbf{380}                                                                                               \\
		Latón y cobre                & 40                                       & 50          & 350          & 400 \\
		Aluminio                     & 60                                       & 90          & 500          & 700 \\
		Fibra y ebonita              & 25                                       & 40
		                             & 120                                      & 150                              \\
	\end{tabular}
\end{table}

\subsection*{Anexo B. Nomenclatura}

\begin{table}[!htb]
	\centering
	\small
	\rowcolors{2}{white}{lightblue}
	\begin{tabularx}{\linewidth}{c X c}
		\rowcolor{darkblue}
		\encab{Símbolo}       & \encab{Significado}                                              & \encab{Unidad} \\
		$v$                   & Velocidad de corte
		                      &
		m/min                                                                                                     \\
		$n$                   & Velocidad de giro del husillo
		                      & rpm                                                                               \\
		$D_i,\ D_f$           & Diámetro inicial y final de la pasada                            &
		mm                                                                                                        \\
		$a$                   & Avance por revolución                                            &
		mm/rev                                                                                                    \\
		$v_L$                 & Velocidad lineal o de avance                                     &
		mm/min                                                                                                    \\
		$Q$                   & Caudal de viruta                                                 & cm$^3$/min     \\
		$A$                   & Área de la corona removida                                       &
		mm$^2$                                                                                                    \\
		$p_T$                 & Profundidad total radial a quitar                                & mm             \\
		$N_D,\ N_A$           & Número de pasadas de desbaste y de acabado                       & ---            \\
		$P_D,\ P_A$           & Profundidad de cada pasada de desbaste y de acabado              & mm             \\
		$L$                   & Longitud mecanizada                                              & mm             \\
		$y$                   & Retiro (aproximación) de la herramienta                          & mm             \\
		$\omega,\ \Delta$     & Ángulo de posición y distancia adicional ($\tan\omega=p/\Delta$) & ($^\circ$), mm \\
		$t_p,\ t_m$           & Tiempo de pasada y tiempo de maquinado                           & min            \\
		$r_\varepsilon$       & Radio de punta de la herramienta                                 & mm             \\
		$J_{\max},\ J_{\min}$ & Juego máximo y mínimo del ajuste                                 & \textmu m      \\
	\end{tabularx}
\end{table}

\subsection*{Anexo C. Referencias}

\begin{enumerate}[label={[\arabic*]},ref=\arabic*,leftmargin=2.2em,itemsep=5pt]
	\item\label{bib:groover} Groover, M. P. (2007). \emph{Fundamentos de manufactura moderna:
		      materiales, procesos y sistemas} (3\textdegree ed.). McGraw-Hill.
	\item\label{bib:kalpakjian} Kalpakjian, S., \& Schmid, S. R. (2014). \emph{Manufactura, ingeniería y
		      tecnología} (7\textdegree ed.). Pearson Educación.
	\item\label{bib:iso286} International Organization for Standardization. (2010). \emph{ISO 286-1:2010
		      Geometrical product specifications (GPS) --- ISO code system for tolerances on linear sizes
		      --- Part 1: Basis of tolerances, deviations and fits}. ISO.
	\item\label{bib:iso2768} International Organization for Standardization. (1989). \emph{ISO 2768-1:1989
		      General tolerances --- Part 1: Tolerances for linear and angular dimensions without
		      individual tolerance indications}. ISO.
	\item\label{bib:astm} ASTM International. \emph{ASTM B505/B505M: Standard Specification for
		      Copper-Base Alloy Continuous Castings} (alloy UNS C93200).
	\item\label{bib:oberg} Oberg, E., Jones, F. D., Horton, H. L., \& Ryffel, H. H. (2020).
	      \emph{Machinery's Handbook} (31\textdegree ed.). Industrial Press.
	\item\label{bib:gum} Joint Committee for Guides in Metrology. (2008). \emph{Evaluation of
		      measurement data --- Guide to the expression of uncertainty in measurement} (JCGM 100:2008).
	      BIPM.
	\item\label{bib:clase} Universidad Nacional Mayor de San Marcos. \emph{Semana 4: Procesos de
		      remoción de torno}. Material de clase del curso Procesos de Manufactura.
\end{enumerate}

\end{document}
